\section*{Hoofdstuk \ref{ch:g9:metals-acids-corrosion} --- Metalen: reacties met zuren en corrosie}

\begin{solution}{exo:g9:metals-acids-corrosion:1}
Waterstof: een brandende spaander bij de opening van de buis geeft een
blaffend plofje.
\end{solution}

\begin{solution}{exo:g9:metals-acids-corrosion:2}
IJzer, zink en aluminium reageren; koper en goud niet.
\end{solution}

\begin{solution}{exo:g9:metals-acids-corrosion:3}
Zowel zuurstof als water.
\end{solution}

\begin{solution}{exo:g9:metals-acids-corrosion:4}
Een \omterm{def:g9:ions:ion}{ion} dat tijdens een reactie aanwezig is maar er niet aan meedoet.
Met zoutzuur: het \omterm{def:g9:ions:ion}{chloride-ion} \ce{Cl-}.
\end{solution}

\begin{solution}{exo:g9:metals-acids-corrosion:5}
\ce{Mg(s) + 2H+(aq) -> Mg^{2+}(aq) + H2(g)}; lading $+2$ aan elke kant.
\end{solution}

\begin{solution}{exo:g9:metals-acids-corrosion:6}
Voeg natriumhydroxide toe: een groene \omterm{def:g9:ions:precipitate}{neerslag} van \ce{Fe(OH)2} wijst op
ijzer(II)\omterm{def:g9:ions:ion}{ionen}.
\end{solution}

\begin{solution}{exo:g9:metals-acids-corrosion:7}
Door te koken verdwijnt de \omterm{def:g4:air-a-mixture-of-gases:air}{lucht} die in het water opgelost is, en de
olie voorkomt dat er weer \omterm{def:g4:air-a-mixture-of-gases:air}{lucht} in \omterm{def:g2:mixing-and-dissolving:dissolve}{oplost}: de spijker heeft water, maar
geen zuurstof.
\end{solution}

\begin{solution}{exo:g9:metals-acids-corrosion:8}
Zout dat opgelost is in het water dat tegen de auto spat, laat het
\omterm{def:g9:metals-acids-corrosion:corrosion}{roesten} sneller gaan.
\end{solution}

\begin{solution}{exo:g9:metals-acids-corrosion:9}
Het oxidelaagje dat ontstaat, hecht aan het metaal en sluit het af van
de \omterm{def:g4:air-a-mixture-of-gases:air}{lucht}: de aantasting stopt aan het oppervlak.
\end{solution}

\begin{solution}{exo:g9:metals-acids-corrosion:10}
Het verven, het oliën of invetten, het \omterm{def:g9:metals-acids-corrosion:galvanising}{verzinken} (of het van roestvast
staal maken).
\end{solution}

\begin{solution}{exo:g9:metals-acids-corrosion:11}
\ce{2Al(s) + 6H+(aq) -> 2Al^{3+}(aq) + 3H2(g)}: aan elke kant 2 Al en 6 H,
lading $+6$ aan elke kant. Voor 200 \omterm{def:g7:atoms-and-molecules:atom}{Al-atomen}: $200 \times \frac{3}{2} = 300$
waterstofmoleculen.
\end{solution}

\begin{solution}{exo:g9:metals-acids-corrosion:12}
Het zink rond de kras corrodeert in plaats van het staal: zolang er zink
in de buurt is, wordt het staal gespaard.
\end{solution}

\begin{solution}{pb:g9:metals-acids-corrosion:1}
\textbf{1.} \ce{Zn(s) + 2H+(aq) -> Zn^{2+}(aq) + H2(g)}.

\textbf{2.} Waterstof: een blaffend plofje met een brandende spaander.

\textbf{3.} Natriumhydroxide geeft een witte \omterm{def:g9:ions:precipitate}{neerslag} van
\ce{Zn(OH)2}.

\textbf{4.} De \omterm{def:g9:ions:ion}{chloride-ionen} \ce{Cl-}.

\textbf{5.} $125.6 - 113.5 = \qty{12.1}{g}$.

\textbf{6.} Er ontstaat waterstof zolang er zink met het zuur reageert;
als er geen zink meer is, houdt het bruisen op.

\textbf{7.} Het ijzer van het staal zou ook reageren, met waterstof en
\ce{Fe^{2+}}-ionen als gevolg; natriumhydroxide zou dan een groene
\omterm{def:g9:ions:precipitate}{neerslag} geven.

\textbf{8.} Het zink wordt aangetast in plaats van het staal en houdt
water en \omterm{def:g4:air-a-mixture-of-gases:air}{lucht} ervan weg.

\textbf{9.} $2 \times 10.0 \times 10.0 = \qty{200}{cm^2}$.

\textbf{10.} $12.1 \div 7.13 \approx \qty{1.70}{cm^3}$.

\textbf{11.} $1.70 \div 200 = \qty{0.0085}{cm}$.

\textbf{12.} $0.0085 \times 10\,000 = \qty{85}{\micro m}$: de zinklaag
was ongeveer \qty{85}{\micro m} dik.
\end{solution}
